% ===== PRÉAMBULE CANONIQUE — à coller VERBATIM en tête de CHAQUE .tex =====
\documentclass[11pt,a4paper]{article}
\usepackage[utf8]{inputenc}\usepackage[T1]{fontenc}\usepackage{lmodern}
\usepackage[french]{babel}
\usepackage{amsmath,amssymb,mathtools}\usepackage{icomma}
\usepackage{siunitx}\sisetup{locale=FR}  % \qty{}{}, \si{}, \ang{}
\usepackage{xcolor}\usepackage{colortbl}
\usepackage{tikz}\usepackage{pgfplots}\pgfplotsset{compat=1.18}
\usepackage[siunitx,europeanresistors]{circuitikz} % circuits (élec.)
\usepackage[most]{tcolorbox}\usepackage{enumitem}\usepackage{array,booktabs}
\usepackage[a4paper,margin=2.2cm]{geometry}
\usetikzlibrary{arrows.meta,calc,angles,quotes,positioning,patterns,%
  backgrounds,shapes.geometric,decorations.pathmorphing,decorations.markings,intersections}
\definecolor{primary}{RGB}{222,49,99}\definecolor{primarydark}{RGB}{150,22,55}
\definecolor{defcol}{RGB}{0,114,178}\definecolor{methcol}{RGB}{52,92,148}
\definecolor{excol}{RGB}{0,135,90}\definecolor{attncol}{RGB}{200,40,55}
\tcbset{boxbase/.style={enhanced,breakable,boxrule=0.9pt,arc=2.5pt,
  left=3mm,right=3mm,top=2.3mm,bottom=2.3mm,fonttitle=\bfseries,coltitle=white}}
% --- boîtes TUTORIEL (noms EXACTS, ne pas renommer) ---
\newtcolorbox{cle}[1][]{boxbase,colback=defcol!6,colframe=defcol,
  title={\textsc{À retenir}\ifx\relax#1\relax\relax\else\textmd{ : \itshape #1}\fi}}
\newtcolorbox{methode}[1][]{boxbase,colback=methcol!7,colframe=methcol,sharp corners,
  title={\textsc{Méthode}\ifx\relax#1\relax\relax\else\textmd{ --- \itshape #1}\fi}}
\newtcolorbox{exemplebox}[1][]{boxbase,colback=excol!5,colframe=excol,
  title={\textsc{Exemple}\ifx\relax#1\relax\relax\else\textmd{ : \itshape #1}\fi}}
\newtcolorbox{piege}[1][]{boxbase,colback=attncol!6,colframe=attncol,
  title={\textsc{Piège}\ifx\relax#1\relax\relax\else\textmd{ : \itshape #1}\fi}}
\newtcolorbox{remarque}[1][]{boxbase,colback=black!4,colframe=black!45,
  title={\textsc{Remarque}\ifx\relax#1\relax\relax\else\textmd{ : \itshape #1}\fi}}
% --- boîtes EXERCICE / CORRIGÉ (noms EXACTS, auto-comptées) ---
\newtcolorbox[auto counter]{exo}[1][]{boxbase,colback=primary!4,colframe=primary,
  title={\textsc{Exercice}~\thetcbcounter\ifx\relax#1\relax\relax\else\textmd{ --- \itshape #1}\fi}}
\newtcolorbox[auto counter]{corrige}[1][]{boxbase,colback=excol!4,colframe=excol,
  title={\textsc{Corrigé}~\thetcbcounter\ifx\relax#1\relax\relax\else\textmd{ --- \itshape #1}\fi}}
\newcommand{\vv}[1]{\vec{#1}}\newcommand{\ds}{\displaystyle}
\newcommand{\R}{\mathbb{R}}\newcommand{\N}{\mathbb{N}}\newcommand{\Z}{\mathbb{Z}}
\begin{document}
\sisetup{per-mode=symbol}

\begin{exo}[deux positions remarquables]
Une lentille convergente a $f=+\qty{15}{\cm}$.
\begin{enumerate}[label=\alph*)]
  \item Un objet est \textbf{à l'infini}. Où se forme l'image ?
  \item Un objet est placé en $2F$ (à $u=+\qty{30}{\cm}$). Calcule $v$ et $\gamma$.
\end{enumerate}
\end{exo}

\begin{exo}[deux façons de trouver $f$]
\begin{enumerate}[label=\alph*)]
  \item Un objet est à \qty{1}{\m} d'une lentille ; l'image, recueillie sur un écran, est à \qty{1}{\m}
        de l'autre côté. Calcule $f$.
  \item Un objet est à \qty{5}{\cm} d'une lentille ; son image se forme à \qty{2.5}{\cm} \emph{du même
        côté} que l'objet. Calcule $f$ et donne la famille de la lentille.
\end{enumerate}
\end{exo}

\begin{exo}[l'image agrandie neuf fois]
Un objet est placé à \qty{1}{\m} d'une lentille convergente. L'image réelle formée de l'autre côté est
\textbf{neuf fois plus grande} que l'objet.
\begin{enumerate}[label=\alph*)]
  \item Quel est le grandissement $\gamma$ (avec son signe) ? En déduire la distance image $v$.
  \item Calcule la distance focale $f$.
\end{enumerate}
\end{exo}

\begin{exo}[la loupe en chiffres]
Une loupe est une lentille convergente de focale $f=+\qty{10}{\cm}$. On place une lettre de hauteur
$o=\qty{2}{\mm}$ à $u=+\qty{6}{\cm}$ de la loupe.
\begin{enumerate}[label=\alph*)]
  \item Calcule $v$. L'image est-elle réelle ou virtuelle ?
  \item Calcule $\gamma$ et la hauteur $i$ perçue. Conclusion ?
\end{enumerate}
\end{exo}

\begin{exo}[une lentille divergente]
Une lentille divergente a $f=-\qty{30}{\cm}$. Un objet de hauteur $o=\qty{4}{\cm}$ est placé à
$u=+\qty{45}{\cm}$.
\begin{enumerate}[label=\alph*)]
  \item Calcule $v$ et précise la nature (réelle/virtuelle) de l'image.
  \item Calcule $\gamma$ et la hauteur $i$. Décris l'image en un mot.
\end{enumerate}
\end{exo}

\end{document}
